%==============================================================================
% Section 8: Conclusion
%==============================================================================
\section{Conclusion}
\label{sec:conclusion}

In this paper we have constructed a Euclidean geometric response dictionary
and the corresponding selection rules for the four Thurston geometries
$\Tthree$, $\Nilthree$, $\Sthree$, and $\Solthree$.  The principal result is
the organisation of the response into a two-layer structure that closes the
bulk inventory and the boundary/cobordism inventory across the four
geometries.

In the first layer we provided a 36-entry bulk inventory and its selection
rules.  Here the Maxwell coefficient
\begin{equation}
  \KMxw = -\frac{L^2}{2r_0^4}
\end{equation}
is common to the four geometries, while the activation conditions for CS,
the Higgsing pattern, the EC slice-minimum structure, and the reduced
$\eta$-mode defect-localisation benchmark branch in a geometry-dependent
manner.  $\Nilthree$ is organised as a CS+EC activated geometry, $\Sthree$
as the most richly CS- and Higgsing-activated geometry, $\Tthree$ as an
inert benchmark, and $\Solthree$ as a CS-inert / EC-active geometry with a
rigid scaffold.

In the second layer we fixed the three observable families spectral,
torsional-charge, and inflow, together with the pairwise cobordism
dictionary.  In particular, the strict local APS core of $\Nilthree$, the
frame-bundle-normalised torsional-charge core of $\Sthree$, the trivial
spectral benchmark of $\Tthree$, and the spin-sensitive global spectral
branch on the compact $\Solthree$ benchmark constitute the centre of the
comparison dictionary.  This makes explicit the separation of local and
global sources within the spectral family.  The geometric scaffold layer
distinguishes this difference of source from the distinctness of the
background geometry.

We have further presented two readings of the Euclidean circle, the KK
reduction reading and the Matsubara-style thermal reading.  These provide a
reduced description and a thermal description respectively, but are not
identified theoretically.  They coexist as secondary interpretations of the
already constructed Euclidean response dictionary.

In conclusion, the present paper offers a Euclidean geometric response
dictionary that gives the allowed, protected, and activated responses of the
four geometries as a coherent rule set.  The principal contribution is the
explicit display of how the layers of bulk, boundary, cobordism, and circle
readings are organised once the four Thurston geometries are placed within
the same comparison framework.
